\section{Graph IR Lowering}

\paragraph{Input and output.}
Graph lowering consumes a flat, closed, fully typed Axon AST and produces an
in-memory typed Graph IR program.

\paragraph{Contract.}
The lowering stage translates Axon definitions to graph modules, bind
statements to graph nodes, multi-output expressions to structured graph
outputs, and path literals to structured graph paths.  It preserves operand
types, output types, constraints, constants, and module boundaries needed by
backends and later graph-level analysis.

\paragraph{Exclusions.}
Lowering does not accept surface constructs such as \code{scope}, \code{for},
or statement-level \code{if}.  It also does not rewrite primitive calls into
legacy string-based callee forms.

\paragraph{Rationale.}
Graph IR replaces the old serialized graph contract.  The compiler should pass
typed structured objects to backends rather than relying on YAML-shaped
dictionaries or render/parse cycles.

\paragraph{Formal view.}
Lowering is a homomorphism from typed flat modules to graph modules:
\[
  \mathsf{lower}: A_{\mathrm{flat}}^{\tau} \rightarrow G.
\]
Each flat definition becomes a graph module.  Each flat bind becomes one graph
node with typed operands and typed outputs.  Calls to user definitions become
module-call nodes; calls to primitive wrappers lower to primitive graph ops
only when the source definition is itself a primitive wrapper.

\paragraph{Example.}
A typed flat bind
\begin{verbatim}
tok <- NN.embedding @@wte input_ids dim=MODEL_DIM
\end{verbatim}
lowers to a graph node whose op is the embedding primitive, whose first operand
is a structured absolute path \code{@@wte}, and whose output type is the
inferred tensor type, e.g. \verb|Tensor[B,S,MODEL_DIM]|.

\paragraph{Multi-output nodes.}
Graph IR preserves multi-output structure.  A flat destructuring bind such as
\begin{verbatim}
q, k, v <- Tensor.chunk qkv parts=3
\end{verbatim}
can lower to a single graph node with three outputs rather than a synthetic
list plus three unrelated index operations.
